\SGANextSectionTarget{sga3:xxiii:subsection:3-5}\SGASubsection{3.5}{Explicit Form of the Generators-and-Relations
Theorem}

\SGARefTarget{sga3:XXIII:3.5.1}
\SGAStatement{Theorem 3.5.1}{Let \(S\) be a scheme and \(G\) a
pinned \(S\)-group, with maximal torus \(T\), system of simple roots
\(\Delta\), and elements
\[
  u_\alpha\in U_\alpha(S)^\times,\qquad
  w_\alpha\in\operatorname{Norm}_G(T)(S)
  \qquad(\alpha\in\Delta)
\]
defined by the pinning. Let \(H\) be an fppf \(S\)-sheaf in groups,
let
\[
  f_T:T\to H,\qquad f_\alpha:U_\alpha\to H
  \quad(\alpha\in\Delta)
\]
be morphisms of groups, and let \(h_\alpha\in H(S)\). Put
\[
  v_\alpha=f_\alpha(u_\alpha).
\]
There exists a morphism of groups \(f:G\to H\) extending \(f_T\) and
all the \(f_\alpha\), and satisfying
\(f(w_\alpha)=h_\alpha\), if and only if:
\begin{enumerate}
  \item[(i)] for every \(S'\to S\), \(\alpha\in\Delta\),
    \(t\in T(S')\), and \(x\in U_\alpha(S')\),
    \begin{equation}
      \operatorname{int}(f_T(t))f_\alpha(x)
       =f_\alpha(\operatorname{int}(t)x)
       =f_\alpha(x^{\alpha(t)}); \tag{1}\label{sga3:xxiii:theorem:3-5-1:equation:1}
    \end{equation}
  \item[(ii)] for every \(\alpha\in\Delta\), every \(S'\to S\), and
    every \(t\in T(S')\),
    \begin{equation}
      \operatorname{int}(h_\alpha)f_T(t)
       =f_T(s_\alpha(t))
       =f_T\!\left(t\alpha^*(\alpha(t)^{-1})\right); \tag{2}\label{sga3:xxiii:theorem:3-5-1:equation:2}
    \end{equation}
  \item[(iii)] for every \(\alpha\in\Delta\),
    \begin{equation}
      h_\alpha^2=f_T(\alpha^*(-1)), \tag{\(3_1\)}\label{sga3:xxiii:theorem:3-5-1:equation:3-1}
    \end{equation}
    and
    \begin{equation}
      (h_\alpha v_\alpha)^3=e; \tag{4}\label{sga3:xxiii:theorem:3-5-1:equation:4}
    \end{equation}
  \item[\SGAInlineRefTarget{sga3:XXIII:3.5.1:item:iv}(iv)] for distinct \(\alpha,\beta\in\Delta\):
    \begin{enumerate}
      \item[(a)] the appropriate braid relation is
        \begin{equation}
        \tag{\(3_2\)}\label{sga3:xxiii:theorem:3-5-1:equation:3-2}
        \begin{cases}
        (h_\alpha h_\beta)^2
          =f_T(\alpha^*(-1)\beta^*(-1)),
          &(\alpha^*,\beta)=0,\\
        (h_\alpha h_\beta)^3=e,
          &(\alpha^*,\beta)=-1,\\
        (h_\alpha h_\beta)^4=f_T(\alpha^*(-1)),
          &(\alpha^*,\beta)=-2,\\
        (h_\alpha h_\beta)^6=e,
          &(\alpha^*,\beta)=-3;
        \end{cases}
        \end{equation}
      \item[(b)] the corresponding relations (A) and (B) of
        \SGARefLink{sga3:XXIII:3.1.3}{3.1.3}, \SGARefLink{sga3:XXIII:3.2.8}{3.2.8}, \SGARefLink{sga3:XXIII:3.3.7}{3.3.7}, or \SGARefLink{sga3:XXIII:3.4.10}{3.4.10} hold.
    \end{enumerate}
\end{enumerate}
When it exists, \(f\) is necessarily unique.}

This follows immediately from \SGARefLink{sga3:XXIII:2.3}{Theorem 2.3}, \SGARefLink{sga3:XXIII:2.6}{Remark 2.6}, and the
calculations in each rank-two case.

\SGARefTarget{sga3:XXIII:3.5.2}
\SGAStatement{Remark 3.5.2}{The result may equivalently be presented
by giving morphisms
\[
  a_\alpha:T\!\cdot U_\alpha\longrightarrow H,
  \qquad
  b_\alpha:\operatorname{Norm}_{Z_\alpha}(T)\longrightarrow H
  \qquad(\alpha\in\Delta),
\]
and putting
\[
  h_\alpha=b_\alpha(w_\alpha),
  \qquad
  v_\alpha=a_\alpha(u_\alpha).
\]
The required conditions are:
\begin{enumerate}
  \item[(1)] all the \(a_\alpha\) and \(b_\alpha\) have the same
    restriction to \(T\);
  \item[(2)] condition (4) and all the conditions in \SGARefLink{sga3:XXIII:3.5.1:item:iv}{(iv)} of
    \SGARefLink{sga3:XXIII:3.5.1}{Theorem 3.5.1} hold.
\end{enumerate}}

\SGARefTarget{sga3:XXIII:3.5.3}
\SGAParagraph{3.5.3}
Further applications will be given in the next Exposé. Here is one.
\SGARefLink{sga3:XXIII:3.5.1}{Theorem 3.5.1} describes \(G\) by generators and relations in the
category of fppf \(S\)-sheaves. More explicitly, for every
\(S'\to S\), let \(H(S')\) be the group generated by
\[
  T(S'),\qquad U_\alpha(S')\ (\alpha\in R),\qquad
  w_\alpha\ (\alpha\in R),
\]
subject to relations analogous to
\[
  (1),\ (2),\ (3_1),\ (4),\ (3_2),\ (A),\ (B).
\]
Then \(G\) is the sheaf associated with the presheaf
\(S'\mapsto H(S')\).

In particular, if \(S'\) is the spectrum of an algebraically closed
field \(k\), then \(G(S')=H(S')\). This follows immediately from the
Nullstellensatz in the form: an fppf-covering sieve over an
algebraically closed field is trivial. Thus \SGARefLink{sga3:XXIII:3.5.1}{Theorem 3.5.1} gives an
explicit presentation by generators and relations of the abstract
group \(G(k)\).
